% closing_remarks.tex - Final chapter closing and resources
% ============================================================================

\section{Where to Go From Here}
\label{sec:next-steps}

This book gave you the essentials. The real learning happens when you start using Python for your own research. Here are resources to help you continue.

\subsection{Packages to Explore}

These libraries extend your toolkit for specific tasks. With what you've learned, you can teach yourself any of them from their documentation:

\begin{itemize}
    \item \textbf{seaborn} --- Statistical visualization built on Matplotlib. Less code, beautiful plots.
    \item \textbf{scikit-learn} --- Machine learning.
    \item \textbf{statsmodels} --- Statistical modeling: linear models, time series, hypothesis tests.
    \item \textbf{sympy} --- Symbolic mathematics: algebra, calculus, equation solving (Mathematica like).
    \item \textbf{lmfit} --- Curve fitting with parameter constraints, confidence intervals, \& model evaluation.
    \item \textbf{uncertainties} --- Automatic error propagation through calculations.
    \item \textbf{pint} --- Physical units handling (never confuse meters and feet again).
\end{itemize}

\subsection{Domain-Specific Libraries}

Search for Python packages in your field---they almost certainly exist:
\begin{itemize}
    \item \textbf{Chemistry:} RDKit, Open Babel, cclib, PyMOL
    \item \textbf{Biology:} Biopython, scikit-bio, scanpy
    \item \textbf{Physics/Astronomy:} Astropy, QuTiP, PyMC
    \item \textbf{Geoscience:} xarray, cartopy, MetPy
    \item \textbf{Image Analysis:} scikit-image, OpenCV, Pillow
\end{itemize}

\subsection{Improve Your Workflow}

As your projects grow, these practices will save you time and headaches:
\begin{itemize}
    \item Learn Git for version control
    \item Write tests for your code
    \item Create reusable packages from your scripts
    \item Explore Jupyter notebooks for interactive analysis
\end{itemize}



\subsection{Practice}

The best way to learn is to use Python for something you actually care about. Pick a dataset from your research. Automate something tedious. Recreate a figure from a paper. Break things, fix them, and learn by doing.

You now have enough foundation to teach yourself the rest.

\subsection{Where to Get Help}
If you do get stuck (and you will be---everyone does) in order of quality:
\begin{itemize}
    \item \textbf{Official Documentation} --- NumPy, Pandas, and Matplotlib all have excellent tutorials and example galleries.
	\item \textbf{Your colleagues}---They will often know more about the specific problems you're trying to solve than most others.
    \item \textbf{Stack Overflow} --- Search before asking. Your error message has almost certainly been answered.
    \item \textbf{GitHub Issues} --- For package-specific bugs or feature questions.
    \item \textbf{Reddit} --- r/learnpython is friendly to beginners; r/Python for general discussion.
	\item \textbf{Chat-bots}---Undeniably useful, even if these can be problematic. Use with caution, and don't accept code you don't understand. These are a double-edged blade, you WILL get cut using these.
\end{itemize}

\begin{warning}
While AI chatbots can be helpful, they are a double-edged sword for beginners. To get a functional result, you often have to describe your problem in such extreme detail that you might as well have written the code yourself. In fact, Python is likely the most efficient language for that description, so with the description done, you have what you wanted. 

Anything less than a perfect prompt usually results in "black box" code: bloated, difficult to debug, and hard to learn from. However, they can be excellent tools if used as a \textbf{reference} rather than a \textbf{ghostwriter}. Chat-bots give you speed, and so if you're bad at the task, now you're just bad faster. Do not mistake speed for efficiency or effectiveness.

\textbf{Chatbots may be effective for:}
\begin{itemize}
    \item \textbf{Search:} ``Is there a NumPy function that calculates the standard deviation?''
    \item \textbf{Syntax:} ``How do I format a f-string to show two decimal places?''
    \item \textbf{Error Decoding:} Pasting an error message to get a plain-English explanation of what it means. \textbf{AFTER} you have tried to understand yourself a few times. Do not default to a bot when running into problems, it robs you of skills and experiences that are worth more than whatever deadline feels important right now.
\end{itemize}

\textbf{Avoid Chatbots for:}
\begin{itemize}
    \item \textbf{Logic:} Asking it to write a whole script. You lose the "mental muscle" built by solving the puzzle yourself. Besides, given what you do, you are quite likely better at logic than a chat-bot.
    \item \textbf{Debugging:} Letting it guess fixes for your code. You will often spend more time trying to fix the chatbot's "fix" than if you had used \texttt{print()} statements to find the bug yourself.
\end{itemize}
\end{warning}

\vfill
\begin{center}
    \textit{Happy Coding!}
\end{center}
